% Fragmento ES para inserción, no documento autónomo.
% Destino: XII, ampliacion_20260922.tex, después de la prueba de
% «Realización exacta del cuarto de giro» y antes de introducir S_zeta.
% Propietarios y comprobaciones: LEER_PRIMERO.md.
\subsubsection{Restitution du calendrier et retenue ternaire}
\label{ins29:xii:acarreo}

L’extraction du plan de quart de tour admet une description exacte
de l’information temporelle conservée. Nous recevons le calendrier de
\(108\) positions et l’extension dodécaphasique sur la phase nonadique
construits précédemment. Pour cette description, nous utilisons des représentants
indexés à partir de zéro :
\[
 \theta=r+9\beta,\qquad 0\leq r<9,\quad 0\leq\beta<12.
\]
La position que le registre énumère de \(1\) à \(9\) est ici \(r+1\).
La mise à jour enrichie précède ce calendrier fini et conserve,
en outre, la route, les incidences et la mémoire accumulée. Dans le plan déjà
construit, \(C_{12}^{\beta}W=WJ_0^{\beta}\) ; sa lecture dépend donc
de \(b=\beta\bmod4\).

\begin{proposition}[Fibre ternaire de la lecture de quart de tour]
\label{ins29:xii:fibra}
Le couple formé par la phase nonadique \(r\) et la phase planaire \(b\) détermine
un quotient de \(36\) positions du calendrier. Sa loi de composition est
\[
 (r,b)(r',b')=
 \left(r+r'\bmod9,\;
 b+b'+\left\lfloor\frac{r+r'}9\right\rfloor\bmod4\right).
\]
Chaque lecture possède trois préimages dans \(\mathbb Z/108\mathbb Z\).
La composante supplémentaire \(k=\beta\bmod3\) restitue la coordonnée
dodécaphasique au moyen de
\[
 \beta=9b+4k\pmod{12}.
\]
\end{proposition}
\begin{proof}
La somme de deux représentants donne
\(\theta+\theta'=r+r'+9(\beta+\beta')\). La division euclidienne de
\(r+r'\) par \(9\) produit la retenue de la formule. La lecture
\((r,b)\) équivaut au résidu \(\theta\bmod36\), puisque
\(\theta\equiv r+9b\pmod{36}\). Son noyau est
\(\{0,36,72\}\), de sorte que chaque fibre a trois éléments.
Enfin, \(9b+4k\) est congru à \(b\) modulo \(4\) et à
\(k\) modulo \(3\) ; ces deux résidus déterminent un unique élément
modulo \(12\).
\end{proof}

\begin{proposition}[Localisation de l’extension non scindée]
\label{ins29:xii:extension}
La décomposition
\[
 \theta\longmapsto(j,z)=(\theta\bmod4,\theta\bmod27),
 \qquad \theta=81j+28z\pmod{108},
\]
identifie le calendrier à
\(\mathbb Z/4\mathbb Z\times\mathbb Z/27\mathbb Z\).
La projection nonadique est \((j,z)\mapsto z\bmod9\), et l’inclusion
dodécaphasique \(\beta\mapsto9\beta\) prend la forme
\[
 \beta\longmapsto(\beta\bmod4,9\beta\bmod27).
\]
L’obstruction à scinder l’extension complète se concentre dans
\[
 0\longrightarrow\mathbb Z/3\mathbb Z
 \longrightarrow\mathbb Z/27\mathbb Z
 \longrightarrow\mathbb Z/9\mathbb Z\longrightarrow0.
\]
\end{proposition}
\begin{proof}
Les coefficients \(81\) et \(28\) ont les résidus \((1,0)\) et \((0,1)\)
modulo \((4,27)\), respectivement. Cela vérifie l’inverse et les deux
applications indiquées. La suite ternaire ne se scinde pas : un produit
\(\mathbb Z/3\mathbb Z\times\mathbb Z/9\mathbb Z\) a pour exposant
\(9\), tandis que \(\mathbb Z/27\mathbb Z\) a pour exposant \(27\).
En revanche, la projection du quotient de \(36\) positions sur
\(\mathbb Z/9\mathbb Z\) admet la section homomorphe
\(n\mapsto28n\pmod{36}\). La réduction élimine donc précisément
cette obstruction ternaire.
\end{proof}

La distinction s’observe déjà dans \(\theta=0,36,72\) : les trois positions
donnent \(r=b=0\), tandis que \(z\) vaut \(0,9,18\), respectivement.
La restitution du calendrier requiert de conserver la fibre que le plan
identifie. En outre, le résidu quaternaire global satisfait
\(j=r+b\pmod4\) ; la phase planaire \(b\) est transportée avec la phase
nonadique et sa retenue.

La portée du quart de tour se trouve ainsi localisée : il organise une
représentation orientée du registre et, accompagné du résidu nonadique,
retrouve exactement son quotient de \(36\) positions. La coordonnée
ternaire supplémentaire restitue les \(108\) positions. La mémoire de la
dynamique enrichie conserve à son tour le prolongement entre les retours
complets du calendrier. La périodicité d’une représentation et la
persistance de la trajectoire enregistrée appartiennent donc à des lectures distinctes
du même transport.
